\documentclass[10pt,a4paper]{amsart}
\usepackage[margin=19mm]{geometry}
\usepackage{amsmath,amssymb,mathtools}
\usepackage[scr=rsfs,cal=euler]{mathalfa}
\usepackage{xfrac}
\usepackage[unicode,hidelinks,bookmarks=false]{hyperref}
\newcommand{\ie}{i.e.\ }
\newcommand{\cf}{cf.\ }
\newcommand{\resp}{respectively\ }
\newcommand{\Subsection}{\subsection}
\providecommand{\nobreakdash}{\nobreak\hbox{-}}
\setlength{\emergencystretch}{2em}
\multlinegap0pt
\usepackage{amssymb}
\usepackage[all]{xy}
\usepackage[shortlabels]{enumitem}
\usepackage{mathtools}
\newtheorem{lem}{Lemma}
\title{On the geometry of holomorphic curves and complex surface}
\author{Amanda Dias Falqueto, Farid Tari}
\date{Source: arXiv:2505.24719v2 (CC BY 4.0)}
\begin{document}
\maketitle

\noindent\emph{Source and licence.} On the geometry of holomorphic curves and complex surface. Version arXiv:2505.24719v2 (CC BY 4.0), distributed by arXiv under the Creative Commons Attribution 4.0 International licence, http://creativecommons.org/licenses/by/4.0/, as recorded in the arXiv licence metadata for this exact version and shown on the abstract page https://arxiv.org/abs/2505.24719v2. Full licence notice: \texttt{licenses/arxiv-2505.24719-CC-BY-4.0.txt}.\par\smallskip

\noindent\textit{Editorial scope.} Counted unit: the article's lem lem:isotPlane (source0.tex lines 584--588). The article's complete original proof of it is delivered with it (source0.tex lines 590--604, one \texttt{proof} environment of 15 lines). No mathematics is written, repaired or re-derived: the statement and the proof are the article's own lines, in the article's own notation, and no retained line is substituted or re-flowed.  No further source-local context is needed: every cross-reference of every retained line resolves inside the two delivered spans.  Nothing was omitted from any retained span, so the package carries no omission record.  No retained line contains a citation, so the wrapper carries no bibliography and no bibliography entry.  No retained line contains a figure environment, an image-inclusion command or drawing code, so no image is delivered and the package declares no asset dependency.  Editorial formatting only, in addition: an AMS \texttt{amsart} preamble that restates the article's own declarations and macros used by the retained lines, namely the environments \texttt{lem} on separate counters, together with the standard packages those lines need.  The retained environments print in this document as Lemma 1 in order of appearance, whereas the article prints the same items with the numbering of its own full document.  The complete native source of the article as distributed in the arXiv e-print for this exact version is bundled unchanged as \texttt{sources/179053/source0.tex}.
\begin{lem}\label{lem:isotPlane}
Suppose that $\gamma$ is parametrised by unit speed at a non-isotropic point $s_0$, and let $T$ be the unit vector of $\gamma$. 
Then the osculating plane $\pi$ at $s_0$ is isotropic 
if and only if $T'(s_0)$ is isotropic.
\end{lem}

\begin{proof}
Assume first that $\pi$ is isotropic. Then $\pi$ has an isotropic normal vector 
$v\nparallel T(s_0)$, so it is generated by $T(s_0)$ and
$v$, with $\langle T(s_0),v\rangle=0$. 

As $T'(s_0)$ is parallel to $\pi$, we can write $T'(s_0)=\alpha T(s_0)+\beta v$ for some $\alpha,\beta\in \mathbb C$.
Differentiating the identity $\langle T,{T}\rangle=1$ gives 
 $\langle T',T\rangle=0$, hence $\alpha=0$. Thus, $T'(s_0)=\beta v$ 
 which is isotropic. 
(We are assuming $\pi$ to be an isotropic plane, so $\beta\ne 0$.)
 
Conversely, suppose that $T'(s_0)$ is an isotropic vector. For any 
$\alpha,\beta\in \mathbb C$, we have $\langle \alpha T(s_0)+\beta T'(s_0),T'(s_0)\rangle=0$.
Thus, $\pi$ is an isotropic plane with normal $T'(s_0)$.
\end{proof}

\end{document}
